%&latex
\documentclass[a4paper]{article}

\frenchspacing

\usepackage[cp1250]{inputenc}
\usepackage[czech]{babel}

\usepackage{a4wide}
\usepackage{amsmath, amsthm, amssymb, amsfonts}
\usepackage[mathcal]{eucal}
\usepackage{graphicx}
\usepackage{url}
\usepackage{color}
\usepackage{wrapfig}
\usepackage{capt-of}
\usepackage{float}
\usepackage{ifthen}
\usepackage{subfig}
\usepackage[normalem]{ulem}


% sirka a vyska textu nastavena jako papir, vsechny okraje vynulovany a pridano 20pt na kazdou stranu
% horizontalni rozmery
\setlength{\textwidth}{\paperwidth}
\addtolength{\textwidth}{-40pt}
\addtolength{\hoffset}{-1in}
\addtolength{\hoffset}{20pt}
\setlength{\oddsidemargin}{0in}
\setlength{\marginparsep}{0in}
% vertikalni rozmery
\setlength{\textheight}{\paperheight}
\addtolength{\textheight}{-60pt}
\addtolength{\voffset}{-1in}
\addtolength{\voffset}{20pt}
\setlength{\topmargin}{0in}
\setlength{\headheight}{0in}
\setlength{\headsep}{0in}


%Obrazek na miste
%pouziti
%%\obrazeknahore{adresa}{popisek}{label}
\long\def\obrazeknahore#1#2#3 {

\begin{figure}[t]
    \centering
    \includegraphics[width=0.8\textwidth]{#1}
    
    \caption{#2}
    \label{#3}
    
\end{figure}

}


%==========================================
%PEKELNA MAKRA NA ZAROVNANI OBRAZKU DOPRAVA

\makeatletter


%tohle je makro, ktere mi dovoluje obtekani i u kratkych environmentu
%ABSOLUTNE nechapu, jak to funguje, ale funguje to
%viz http://tex.stackexchange.com/questions/26078/ 
\def\odrovnej{\@@par
\ifnum\@@parshape=\z@ \let\WF@pspars\@empty \fi % reset `parshape'
\global\advance\c@WF@wrappedlines-\prevgraf \prevgraf\z@
\ifnum\c@WF@wrappedlines<\tw@ \WF@finale \fi}

\makeatother



%---
%makro, co da obrazek doprava a ostatni text ho obteka
%(bez toho predchazejiciho makra to ale poradne nebeha)
%pouziti:
%\obrazekvpravo{adresa}{popisek}{label}{procento sirky}
\long\def\obrazekvpravo#1#2#3#4{

\setlength\intextsep{-20pt}

    \begin{wrapfigure}{r}{#4\textwidth}
      \begin{center}
          \vspace{-10pt}
          
        \includegraphics[width=#4\textwidth]{#1}
        \vspace{-10pt}
        
      \end{center}
      
      \caption{#2}
      \label{#3}
      
      
    \end{wrapfigure}

\setlength\intextsep{0pt}

    
}




%---
%makro pro pripady, kdy wrapfigure neco mrsi
%je to docela pekelne
%je nutne mu dat jak text vpravo, tak text vlevo
%a nevim, jestli bude 100% fungovat, ale doufam, ze jo

%pouziti:
%\obrazekvpravominipage{adresa}{popisek}{label}{procento sirky}{1 - procento sirky}{text vlevo}
\long\def\obrazekvpravominipage#1#2#3#4#5#6{

\noindent\begin{minipage}{#5\linewidth}
\vspace{0pt}
#6
\end{minipage}
\hspace{0.5cm}
\noindent\begin{minipage}{#4\linewidth}
\vspace{0pt}
\centering
\includegraphics[width=0.9\textwidth]{#1}
\captionof{figure}{#2}
\label{#3}
\end{minipage}

}

%KONEC PEKELNYCH MAKER
%=====================

\def\lebka{\includegraphics[width=3mm]{../common/lebka}}

% makra pro poznamku u vyrokove a predikatove logiky
\def\vl{ -- ve v�rokov� logice}
\def\pl{ -- v predik�tov� logice}


%Vacsina prostredi je dvojjazicne. V pripade, ze znenie napr pozorovania je pisane po slovensky, malo by byt po slovensky aj oznacenie.

\def\ifis#1#2{\ifthenelse{\equal{#1}{0}}{}{#2}}

%\newenvironment{e}[3]{\pagebreak[2]\noindent\textbf{\ifis{#2}{$\bigstar$}#1}\ifis{#3}{\emph{~(#3)}}\par\noindent\leftskip 10pt }{\odrovnej\par\bigskip}
\newenvironment{e}[3]{\pagebreak[2]\noindent\textbf{\ifis{#2}{\lebka~}#1}\ifis{#3}{\emph{~(#3)}}\ifis{#2}{~\lebka}\par\noindent\leftskip 10pt }{\odrovnej\par\bigskip}


%obecne prostredie, ktore ma vyuzitie pri specialnych odstavcoch ako (uloha, algoritmus...) aby nevzniklo dalsich x prostredi
\newenvironment{obecne}[1]{\pagebreak[2]\noindent\textbf{#1}\par\noindent\leftskip 10pt}{\odrovnej\par\bigskip}

\newenvironment{report}{\pagebreak[2]\noindent\textbf{Report}\em\par\noindent\leftskip 10pt}{\par\bigskip}

%\newenvironment{reportN}[1]{\pagebreak[2]\noindent\textbf{Report~}\emph{(#1)}\emph\par\noindent\leftskip 10pt}{\odrovnej\par\bigskip}
\newenvironment{reportN}[1]{\pagebreak[2]\noindent\textbf{Report~}\emph{(#1)}\em\par\noindent\leftskip 10pt}{\odrovnej\par\bigskip}

\newenvironment{penumerate}{
\begin{enumerate}
  \setlength{\itemsep}{1pt}
  \setlength{\parskip}{0pt}
  \setlength{\parsep}{0pt}
  %\setlength{\topsep}{200pt}
  \setlength{\partopsep}{200pt}
}{\end{enumerate}}

\def\pismenka{\numberedlistdepth=2} %pouzit, ked clovek chce opismenkovany zoznam...

\newenvironment{pitemize}{
\begin{itemize}
  \setlength{\itemsep}{1pt}
  \setlength{\parskip}{0pt}
  \setlength{\parsep}{0pt}
}{\end{itemize}}

\definecolor{gris}{gray}{0.95}
\newcommand{\ramcek}[2]{\begin{center}\fcolorbox{white}{gris}{\parbox{#1}{#2}}\end{center}\par}


\title{\LARGE U�ebn� texty k st�tn� bakal��sk� zkou�ce \\ Matematika \\ Integr�l}

\begin{document}

\setcounter{section}{3}
\def\d{\mathrm{d}}
\def\vol{\mathrm{vol\ }}


\section{Integr�l}

\begin{e}{Po�adavky}{0}{0}
\begin{pitemize}
        \item Primitivn� funkce, metody v�po�tu
        \item Ur�it� (Riemann�v) integr�l, u�it� ur�it�ho integr�lu
        \item V�cerozm�rn� integr�l a Fubiniho v�ta
\end{pitemize}
\end{e}

\subsection{Primitivn� funkce, metody v�po�tu}

\begin{e}{Definice}{0}{0}
Nech� funkce $f$ je definov�na na otev�en�m intervalu $I$. �ekneme, �e funkce $F$ je \emph{primitivn� funkce k $f$ na $I$}, jestli�e pro ka�d� $x \in I$ existuje $F'(x)$ a plat� $F'(x)=f(x)$.
\end{e}

\begin{e}{V�ta}{0}{Tvar primitivn� funkce}
Nech� $F$ a $G$ jsou dv� primitivn� funkce k funkci $f$ na otev�en�m intervalu $I$. Pak existuje $c \in \mathbb{R}$ tak, �e $F(x)=G(x)+c$ pro ka�d� $x \in I$.
\end{e}

\begin{e}{V�ta}{0}{Linearita primitivn� funkce}
Nech� f m� na otev�en�m intervalu $I$ primitivn� funkci $F$, funkce $g$ m� na na $I$ primitivn� funkci $G$ a $\alpha, \beta \in \mathbb{R}$. Potom funkce $\alpha F + \beta G$ je primitivn� funkc� k $\alpha f + \beta g$ na $I$.

\begin{e}{Pozn�mka}{0}{0}
P�edchoz� tvrzen� �asto zapisujeme (pokud alespo� jedno z ��sel $\alpha, \beta$ je r�zn� od nuly)
$$\int (\alpha f(x) + \beta g(x)) \d x = \alpha \int f(x) \d x + \beta \int g(x) \d x.$$
\end{e}
\end{e}

\begin{e}{V�ta}{0}{Spojitost a existence primitivn� funkce}
Nech� $f$ je spojit� funkce na otev�en�m intervalu $I$. Pak $f$ m� na $I$ primitivn� funkci.
\end{e}

\begin{e}{V�ta}{0}{O substituci}
\begin{penumerate}
\item Nech� $F$ je primitivn� funkce k $f$ na $(a,b)$. Nech� $\varphi$ je funkce definov�na na $(\alpha, \beta)$ s hodnotami v intervalu $(a,b)$, kter� m� v ka�d�m bod� $t�\in (\alpha, \beta)$ vlastn� derivaci. Pak
$$\int f(\varphi(t))\varphi'(t)dt=F(\varphi(t))+C \textit{ na } (\alpha,\beta)$$
\item Nech� funkce $\varphi$ m� v ka�d�m bod� intervalu $(\alpha,\beta)$ nenulovou vlastn� derivaci a $\varphi((\alpha,\beta)) = (a,b)$. Nech� funkce $f$ je definov�na na intervalu $(a,b)$ a plat�
$$\int f(\varphi(t))\varphi'(t)dt=G(t)+C \textit{ na } (\alpha,\beta)$$
Pak
$$\int f(x)\d x=G(\varphi^{-1}(x))+C \textit{ na } (a,b)$$
\end{penumerate}
\end{e}

\begin{e}{V�ta}{0}{Integrace per partes}
Nech� $I$ je otev�en� interval a funkce $f$ a $g$ jsou spojit� na $I$. Nech� $F$ je primitivn� funkce k $f$ na $I$ a $G$ je primitivn� funkce ke $g$ na $I$. Pak plat�
$$\int g(x)F(x)\d x=G(x)F(x) - \int G(x)f(x)\d x \textit{ na } I$$

(\emph{Pozn�mka autora}: $\int u'v = uv - \int uv'$)
\end{e}

\subsubsection{Postup integrace racion�ln� funkce}

\begin{e}{V�ta}{0}{D�len� polynomu}
Nech� P a Q jsou polynomy s re�ln�mi koeficienty, p�i�em� Q nen� identicky roven nule. Pak existuj� (jednozna�n� ur�en�) polynomy R a S takov�, �e stupe� S je men�� ne� stupe� Q a pro v�echna $x \in \mathbb{R}$ plat� $P(x)=R(x)Q(x)+S(x)$.
\end{e}

\begin{e}{V�ta}{0}{Z�kladn� v�ta algebry}
Nech� $P(x)=a_nx^n+\dots+a_1x+a_0$ je polynom stupn� $n$ s re�ln�mi koeficienty. Pak existuj� ��sla $x_1,\dots,x_n�\in \mathbb{C}$ takov�, �e
$$P(x)=a_n(x-x_1)\dots(x-x_n), \,\, x \in \mathbb{R}$$
\end{e}

\begin{e}{V�ta}{0}{O ko�enech polynomu}
Nech� $P$ je polynom s re�ln�mi koeficienty a $z \in \mathbb{C}$ je ko�en $P$ n�sobnosti $k \in \mathbb{N}$. Pak i $\overline{z}$ je ko�en $P$ n�sobnosti k.
\end{e}

\begin{e}{V�ta}{0}{O rozkladu polynomu}
Nech� $P(x)=a_nx^n +\dots+a_1x+a_0$ je polynom stupn� $n$ s re�ln�mi koeficienty. Pak existuj� re�ln� ��sla $x_1,\dots,x_k,\alpha_1,\dots,\alpha_l,\beta_1,\dots,\beta_l$ a p�irozen� ��sla $p_1,\dots,p_k,q_1,\dots,q_l$ takov�, �e:

\begin{penumerate}
        \item $P(x)=a_n(x-x_1)^{p_1} \dots (x-x_k)^{p_k}(x^2+\alpha_1x+\beta_1)^{q_1}\dots(x^2+\alpha_lx+\beta_l)^{q_l}$,
        \item ��dn� dva z mnoho�len� $x-x_1,x-x_2,\dots,x-x_k,x^2+\alpha_1x+\beta_1,\dots,x^2+\alpha_lx+\beta_l$ nemaj� spole�n� ko�en,
        \item mnoho�leny $x^2+\alpha_1x+\beta_1,\dots,x^2+\alpha_lx+\beta_l$ nemaj� ��dn� re�ln� ko�en.
\end{penumerate}
\end{e}

\begin{e}{V�ta}{0}{O rozkladu na parci�ln� zlomky}
Nech� $P,Q$ jsou polynomy s re�ln�mi koeficienty takov�, �e

\begin{penumerate}
        \item stupe� $P$ je ost�e men�� ne� stupe� $Q$,
        \item $Q(x)=a_n(x-x_1)^{p_1}\dots(x-x_k)^{p_k} (x^2+\alpha_1x+\beta_1)^{q_1}\dots(x^2+\alpha_lx+\beta_l)^{q_l}$,
        \item $a_n,x_1,\dots,x_k,\alpha_1,\dots,\alpha_l,\beta_1,\dots,\beta_l \in \mathbb{R}, a_n \neq 0$
        \item $p_1,\dots,p_k,q_1,\dots,q_l \in \mathbb{N}$
        \item ��dn� dva z mnoho�len� $x-x_1, x-x_2,\dots,x-x_k,x^2+\alpha_1x+\beta_1,\dots,x^2+\alpha_lx+\beta_l$ nemaj� spole�n� ko�en
        \item mnoho�leny $x^2+\alpha_1x+\beta_1, \dots, x^2+\alpha_lx+\beta_l$ nemaj� re�ln� ko�en
\end{penumerate}

Pak existuj� jednozna�n� ur�en� ��sla $A_1^1,\dots,A_{p_1}^{1},\dots,A_1^k,\dots,A_{p_k}^k,\\B_1^1,C_1^1,\dots,B_{q_1}^1,C_{q_1}^{1},\dots,B_1^l,C_1^l,\dots,B_{q_l}^l,C_{q_l}^l$ takov�, �e plat�


\begin{align*}
\frac{P(x)}{Q(x)} & = \frac{A_1^1}{(x-x_1)^{p_1}} + \dots + \frac{A_{p_1}^1}{(x-x_1)} + \dots + \frac{A_1^k}{(x-x_k)^{p_k}} + \dots + \frac{A_{p_k}^k}{(x-x_k)}\\
&+ \frac{B_1^1x+C_1^1}{(x^2+\alpha_1x+\beta_1)^{q_1}} + \dots + \frac{B_{q_1}^1x+C_{q_1}^1}{x^2+\alpha_1x+\beta_1} + \dots\\
&+ \frac{B_1^lx+C_1^l}{(x^2+\alpha_lx+\beta_l)^{q_l}} + \dots + \frac{B_{q_l}^lx+C_{q_l}^l}{x^2+\alpha_lx+\beta_l}.
\end{align*}

\end{e}

\textbf{Postup integrace racion�ln� funkce $\frac{P(x)}{Q(x)}$ je:}
\begin{penumerate}
\item Vyd�l�me polynomy $P$�a $Q$ - najdeme polynomy $R$ a $S$ takov�, �e
        $$\frac{P(x)}{Q(x)} = R(x) + \frac{S(x)}{Q(x)}$$
        a stupe� S je men�� ne� stupe� Q.

\item Najdeme rozklad polynomu Q ve tvaru uveden�m ve \emph{v�t� o rozkladu polynomu}.

\item Najdeme rozklad $\frac{S}{Q}$ na parci�ln� zlomky ve tvaru uveden�m ve \emph{v�t� o rozkladu na parci�ln� zlomky}.

\item Najdeme primitivn� funkce ke v�em parci�ln�m zlomk�m.
\end{penumerate}

\subsection{Ur�it� (Riemann�v) integr�l, u�it� ur�it�ho integr�lu}

\begin{e}{Definice}{0}{0}
Kone�nou posloupnost $D=\{x_j\}_{j=0}^n$ naz�v�me \emph{d�len�m intervalu} $\left<a,b\right>$, jestli�e plat�
$$a=x_0 < x_1 < \dots < x_n = b$$

Body $x_0,\dots,x_n$ naz�v�me d�l�c�mi body. \emph{Normou d�len�} $D$ rozum�me ��slo
$$\upsilon(D)=\max\{x_j-x_{j-1}; j=1,\dots,n\}.$$

�ekneme, �e d�len� $D'$ intervalu $\left<a,b\right>$ je \emph{zjemn�n�m d�len�} $D$ intervalu $\left<a,b\right>$, jestli�e ka�d� d�l�c� bod $D$ je i d�l�c�m bodem $D'$.
\end{e}

\begin{e}{Definice}{0}{0}
Nech� $f$ je omezen� funkce definovan� na intervalu $\left<a,b\right>$ a $D=\{x_j\}_{j=0}^n$ je d�len� $\left<a,b\right>$. Ozna�me
$$S(f,D)=\sum_{j=1}^{n} M_j(x_j-x_{j-1}), \textit{ kde } M_j=\sup\{f(x); x\in \left<x_{j-1},x_j\right>\}$$
$$s(f,D)=\sum_{j=1}^{n} m_j(x_j-x_{j-1}), \textit{ kde } m_j=\inf\{f(x); x\in \left<x_{j-1},x_j\right>\}$$

\begin{e}{Pozn�mka}{0}{0}
Nech� $f$ je omezen� funkce definovan� na intervalu $\left<a,b\right>$.
\begin{penumerate}
        \item Pro ka�d� d�lan� $D$ intervalu $\left<a,b\right>$. plat� $s(f,D)\le S(f,D)$.
        \item Je-li $D_1$ zjemn�n�m $D_2$, pak $s(f,D_1) \ge s(f,D_2)$ a $S(f,D_1) \le S(f, D_2)$
        \item Jsou-li $D_1$ a $D_2$ d�len� intervalu $\left<a,b\right>$, pak $s(f,D_1) \le S(f, D_2)$.
\end{penumerate}
\end{e}
\end{e}

\begin{e}{Definice}{0}{0}
Nech� $f$ je omezen� funkce definovan� na intervalu $\left<a,b\right>$.
\begin{penumerate}
\item
Ozna�me
$$\overline{\int_a^b}f=\inf\{S(f,D); D \textit{ je d�len�m intervalu } \left<a,b\right>\}$$
\begin{center}(tzv. horn� \emph{Riemann�v integr�l funkce} $f$ p�es $\left<a,b\right>$),\end{center}

$$\underline{\int_a^b}f=\sup\{s(f,D); D \textit{ je d�len�m intervalu } \left<a,b\right>\}$$
\begin{center}(tzv. doln� \emph{Riemann�v integr�l funkce} $f$ p�es $\left<a,b\right>$)\end{center}

\item �ekneme, �e funkce $f$ m� \emph{Riemann�v integr�l p�es} $\left<a,b\right>$, pokud $\overline{\int_a^b}f=\underline{\int_a^b}f$. Hodnota tohoto integr�lu je pak rovna $\overline{\int_a^bf}$ a zna��me ji $\int_a^bf$. 

Pokud $a>b$, definujeme $\int_a^bf=-\int_b^af$, v p��pad�, �e $a=b$, definujeme $\int_a^b=0$.
\end{penumerate}
\end{e}

\begin{e}{V�ta}{0}{Krit�rium existence Riemannova integr�lu}
Nech� $a<b$ a $f$ je funkce omezen� na $\left<a,b\right>$. Pak $\int_a^bf$ existuje, pr�v� kdy� pro ka�d� $\varepsilon > 0$ existuje d�len� $D$ intervalu $\left<a,b\right>$ takov�, �e $S(f,D)-s(f,D)<\varepsilon$
\end{e}

\begin{e}{V�ta}{0}{Monotonie a linearita Riemannova integr�lu}
Nech� $a,b\in \mathbb{R}, a<b$ a funkce $f,g$ maj� Riemann�v integr�l p�es interval $\left<a,b\right>$. Pak plat�:
\begin{penumerate}
        \item Jestli�e pro ka�d� $x \in \left<a,b\right>$ je $f(x) \le g(x)$, pak $\int_a^bf�\le \int_a^bg$
        \item $\int_a^b(f+g) = \int_a^bf + \int_a^bg$
        \item $\int_a^bcf=c\int_a^bf$ pro ka�d� $c \in \mathbb{R}$
\end{penumerate}
\end{e}

\begin{e}{V�ta}{0}{Spojitost a Riemannovsk� integrovatelnost}
Nech� $a,b\in \mathbb{R}, a<b$ a funkce $f$ je spojit� na intervalu $\left<a,b\right>$. Pak existuje $\int_a^bf$.
\begin{e}{Pozn�mka}{0}{0}
Plat� dokonce: Pokud je $f$ omezen� na $\left<a,b\right>$ a je spojit� ve v�ech bodech intervalu $\left<a,b\right>$ s v�jimkou kone�n� mnoha, pak existuje $\int_a^bf$.
\end{e}
\end{e}


\begin{e}{V�ta}{0}{Monotonie a Riemannovsk� integrovatelnost}
Je li $f$ omezen� a monot�nn� na uzav�en�m intervalu, pak je Riemannovsky integrovateln�. % Proc ty otazniky? to plati. -- Tuetschek
\end{e}

\begin{e}{V�ta}{0}{Vlastnosti $\int$}
Nech� $a,b\in \mathbb{R}, a<b$ a funkce $f$ je omezen� na intervalu $\left<a,b\right>$. Pak plat�
\begin{penumerate}
        \item Jestli�e existuje $\int_a^bf$, pak pro ka�d� interval $\left<c,d\right> \subset \left<a,b\right>$ existuje $\int_c^df$.
        \item Je-li $c \in (a,b)$, pak $\int_a^bf=\int_a^cf+\int_c^bf$, m�-li alespo� jedna strana smysl (\emph{aditivita Riemannova integr�lu jako funkce intervalu})
\end{penumerate}
\end{e}

\begin{e}{V�ta}{0}{Riemann�v integr�l jako primitivn� funkce}
Nech� $a,b\in \mathbb{R}, a<b$ a funkce $f$ je omezen� na intervalu $\left<a,b\right>$. Pro $x \in \left<a,b\right>$ polo�me $F(x)=\overline{\int_a^x}f$. Potom plat�:
\begin{penumerate}
        \item Funkce $F$ je spojit� na $\left<a,b\right>$
        \item Je-li $x_0 \in (a,b)$ a funkce $f$ je v bod� $x_0$ spojit�, pak $F'(x_0)=f(x_0)$.
\end{penumerate}
Stejn� tvrzen� plat� i pro funkci $G(x)=\underline{\int_a^x}f$
\end{e}

\begin{e}{Pozn�mka}{0}{0}
Pro Riemannovy integr�ly lze pou��t i metodu per partes nebo pravidlo substituce. % tohle by mozna chtelo formulovat nejak formalneji -- Tuetschek
\end{e}


\begin{e}{V�ta}{0}{Z�kladn� v�ta anal�zy}
Nech� $a,b\in \mathbb{R}, a<b$ a funkce $f$ je spojit� na intervalu $\left<a,b\right>$. Nech� $F$ je primitivn� funkce k $f$ na intervalu $(a,b)$. Pak existuj� vlastn� limity $\lim_{x\rightarrow a+}F(x)$ a $\lim_{x\rightarrow b-}F(x)$ a plat�:
$$\int_a^bf = (\lim_{x\rightarrow b-}F(x))-(\lim_{x\rightarrow a+}F(x))$$
\end{e}


\begin{e}{Definice}{0}{Newton�v integr�l}
Nech� funkce $f$ je definov�na na intervalu $(a,b)$ a $F$ je primitivn� funkce k $f$ na $(a,b)$. \emph{Newtonov�m integr�lem} funkce $f$ p�es interval $(a,b)$ naz�v�me ��slo
$$(N)\int_a^bf = (\lim_{x\rightarrow b-}F(x))-(\lim_{x\rightarrow a+}F(x))$$
pokud ob� limity na prav� stran� existuj� a jsou vlastn�.
\end{e}

\begin{e}{Pozn�mka}{0}{Vztah Riemannova a Newtonova integr�lu}
Je-li funkce $f$ spojit� na intervalu $\left<a,b\right>$, pak plat�: 
$$(N)\int_a^b f(x) \d x = (R)\int_a^b f(x) \d x$$

Mno�iny funkc� integrovateln�ch newtonovsky a riemannovsky jsou neporovnateln�.
\end{e}


\subsubsection{U�it� ur�it�ho integr�lu}

Obsahy rovinn�ch �tvar�...

\begin{e}{V�ta}{0}{D�lka k�ivky}
Nech� $f$ m� na $(a,b)$ spojitou derivaci. D�lka k�ivky v $\mathbb{R}^2$, vyzna�en� pr�b�hem funkce $f$ z $[a;f(a)]$ do $[b;f(b)]$ potom je d�na p�edpisem:
$$L(f)=\int_a^b \sqrt{1+(f'(x))^2}\d x.$$
\end{e}

\begin{e}{V�ta}{0}{Objem rota�n�ho t�lesa}
Nech� $f$ je definov�na na $\left<a,b\right>$ a $f>0$. Objem t�lesa vzniknut�ho rotac� k�ivky je $V=\pi \int_a^bf(x)^2\d x$
\end{e}

\begin{e}{V�ta}{0}{Integr�ln� krit�rium konvergence �ad}
Nech� $f$ je spojit�, nez�porn� a nerostouc� na $\left<n_0-1,\infty\right)$, kde $n_0 \in \mathbb{N}$. Potom

$$\sum_{n=1}^{\infty} f(n) \textit{ konverguje} \Leftrightarrow (N)\int_{n_0}^{\infty}f(t)dt < \infty$$
\end{e}

\subsection{V�cerozm�rn� integr�l a Fubiniho v�ta}

\begin{e}{Definice}{0}{0}
(Kompaktn�m) \emph{intervalem} v $n$-rozm�rn�m euklidovsk�m prostoru $E_n$ rozum�me sou�in
$$J = \left<a_1,b_1\right> \times \dots \times \left<a_n,b_n\right>$$
kde $\left<a_i,b_i\right>$ jsou kompaktn� intervaly v $\mathbb{R}$.
\end{e}

\begin{e}{Definice}{0}{0}
\begin{pitemize}
\item \emph{Rozd�len�m} $D$ takov�ho intervalu $J$ rozum�me $n$-tici $D_1, \dots, D_n$, kde $D_i$ je rozd�len� intervalu $\left<a_i,b_i\right>$.
\item Rozd�len� $D=(D_1,\dots,D_n)$ je \emph{zjemn�n�m} $D'=(D'_1,\dots,D'_n)$ jestli�e $D_i$ zjem�uje $D'_i$.
\end{pitemize}
\end{e}

\begin{e}{Pozorov�n�}{0}{0}
Ka�d� dv� rozd�len� maj� spole�n� zjemn�n�.
\end{e}

\begin{e}{Definice}{0}{0}
\emph{�len rozd�len�} $D=(D_1,\dots,D_n)$ je kter�koliv interval $K=\left<t_{1, i_1}, t_{1, i_1+1}\right> \times \dots \times \left<t_{n, i_n}, t_{n, i_n+1}\right>$, kde $D_k: t_{k0} < \dots < t_{k,r(k)}, \, \, 0 \le i_j \le r(j)$. Mno�ina v�ech �len� rozd�len� $D$ bude ozna�ov�na $|D|$. % urcite tu melo byt podrozdeleni?? -- Tuetschek
\end{e}

\begin{e}{Definice}{0}{0}
\emph{Objem intervalu} $J=\left<a_1,b_1\right> \times \dots \times \left<a_n,b_n\right>$ je ��slo
$$\vol J=(b_1-a_1).(b_2-a_2)\dots(b_n-a_n)$$
\end{e}

\begin{e}{Definice}{0}{0}
Bu� $f$ omezen� funkce na intervalu $J$, bu� $D$ rozd�len� $J$. \emph{Doln� (resp. horn�) sumou} funkce $f$ v rozd�len� $D$ rozum�me ��slo
$$s(f,D)=\sum_{K\in |D|}m_K\cdot\vol K \textit{\;\;resp.\;\;} S(f,D) = \sum_{K \in |D|} M_K\cdot\vol K,$$
kde $m_K$ je infimum a $M_k$ supremum funkce $f$ na intervalu $K$.
\end{e}

\begin{e}{Pozorov�n�}{0}{0}
Pro libovoln� dv� rozd�len� $D$ a $D'$ plat� $s(f,D) \le S(f,D')$
\end{e}

\begin{e}{Definice}{0}{0}
Doln� a horn� Riemann�v integr�l definujeme jako
$$\underline{\int}_Jf=\sup_Ds(f,D), \;\;\; \overline{\int}_Jf=\inf_DS(f,D)$$
a p�i rovnosti t�chto hodnot mluv�me o Riemannov� integr�lu a p�eme prost�
$$\int_Jf \textit{ nebo } \int_Jf(x_1,\dots,x_n)\d x_1\dots x_n, \;\;\;\int_Jf(\overrightarrow{x})d\overrightarrow{x}.$$
\end{e}

\begin{e}{V�ta}{0}{0}
Pro v�cerozm�rn� Riemann�v integr�l plat� (podobn� jako pro jednorozm�rn� p��pad) �e $f$ je Riemannovsky integrovateln�, pr�v� kdy� ke ka�d�mu $\varepsilon > 0$ existuje rozd�len� $D$ takov�, �e
$$S(f,D)-s(f,D)<\varepsilon$$

Plat� i v�ta, �e spojit� funkce na intervalu $J$ je Riemannovsky integrovateln�.
\end{e}

\begin{e}{V�ta}{0}{Vlastnosti Riemannova v�cerozm�rn�ho integr�lu}
Plat�:
\begin{penumerate}
        \item $|\int_J f| \le \int_J|f|$ (existuj�-li p��slu�n� integr�ly)
        \item Bu�te $f,g$ Riemannovsky integrovateln� funkce na $J$, bu� $f \le g$. Potom $\int_J f \le \int_J g$.
        \item Speci�ln�, je-li $f(\overrightarrow{x})\le C$ pro n�jakou konstantu $C$, plat� $\int_J f \le C\cdot \vol J$
\end{penumerate}
\end{e}

\begin{e}{V�ta}{0}{Fubiniova}
Bu�te $J' \subseteq E^n, J'' \subseteq E^m$ intervaly, $J=J' \times J''$, bu� $f$ spojit� funkce definovan� na $J$. Potom
$$\int_J f(\overrightarrow x,\overrightarrow y)d\overrightarrow x\overrightarrow y =
\int_{J'}\left(\int_{J''}f(\overrightarrow x,\overrightarrow y)d\overrightarrow y\right)d\overrightarrow x =
\int_{J''}\left(\int_{J'}f(\overrightarrow x,\overrightarrow y)d\overrightarrow x\right)d\overrightarrow y
$$

(In�mi slovami: Hodnota \uv{integr�lu} cez cel� interval je rovn� hodnote po integrovan� postupne cez (jednotliv�) \uv{rozmery} - pri�om je mo�n� integrova� v �ubovo�nom porad�.)
\end{e}

\begin{e}{Definice}{0}{Parci�ln� derivace}
\emph{Parci�ln� derivace} funkce $f$ v bod� $a\in\mathbb{R}^n$ podle prom�nn� $x_i$ se definuje n�sledovn�:
$$\frac{\partial f}{\partial x_i}(a)=\lim_{h\to 0}\frac{f(a_1,\dots,a_{i-1},a_i + h,a_{i+1},\dots,a_n)-f(a_1,\dots,a_n)}{h}$$
\end{e}

\begin{e}{Definice}{0}{Jacobiho matice, jakobi�n}
\emph{Jacobiho matice} funkce $\overrightarrow f:D\to \mathbb{R}^n$ v bod� $a\in D$, kde $D$ je otev�en� mno�ina v $\mathbb{R}^m$ a $f_1,f_2,\dots f_n$ jsou sou�adnicov� funkce $f$, je d�na p�edpisem:
$$
\left(\frac{\partial f_i}{\partial x_j}(a)\right)_{i,j=1}^{n,m} =
\left( \begin{matrix}
        \frac{\partial f_1}{\partial x_1} & \cdots & \frac{\partial f_1}{\partial x_m} \\
        \vdots & \ddots & \vdots \\
        \frac{\partial f_n}{\partial x_1} & \cdots & \frac{\partial f_n}{\partial x_m}
\end{matrix} \right)
$$

\noindent T�to �tvercov� matice se obvykle zna�� $\frac{D(f_1,\dots,f_n)}{D(x_1,\dots,x_n)}$ -- a je-li $m=n$, jej� determinant se naz�v� \emph{Jakobi�n} (a zna�� se rovnako???).
\end{e}

\begin{e}{Definice}{0}{Regul�rn� zobrazen�}
Nech� $U \subseteq \mathbb{R}^n$ je otev�en� mno�ina, $\overrightarrow f: U \rightarrow \mathbb{R}^n$ m� spojit� parci�ln� derivace. \emph{Zobrazen� $\overrightarrow f$ je regul�rn�}, je-li jakobi�n
$$\frac{D(f_1,\dots,f_n)}{D(x_1,\dots,x_n)}(\overrightarrow x) \neq 0, \,\, \forall \overrightarrow x \in U$$
\end{e}


\begin{e}{V�ta}{0}{O substituci}
Nech� $\varphi: U \subseteq \mathbb{R}^n \rightarrow \mathbb{R}^n$ je regul�rn� zobrazen�, $A$ je uzav�en� mno�ina v $\mathbb{R}^n$, $A \subseteq U$ na kter� existuje $\int_{\varphi(A)}f(\overrightarrow x)d\overrightarrow x$. Potom plat�:
$$\int_A f(\overrightarrow{\varphi}(\overrightarrow t)) \frac{D(\overrightarrow \varphi)}{D(\overrightarrow t)}d\overrightarrow t = \int_{\varphi (A)} f(\overrightarrow x)d\overrightarrow x$$
\end{e}

\newpage
\begin{e}{Report}{0}{Klazar} Primitivn� funkce - napsal jsem definici, �e je to line�rn� zobrazen�, per partes a uvedl jsem Z�kladn� v�tu anal�zy. Zkou�ej�c� se je�t� zeptali na v�tu o substituci, tu jsem neznal, co� z�ejm� nijak nevadilo.
\end{e}
\begin{e}{Report}{0}{Majerech} primitivni fce (jen co se tyce zavedeni a pocitani s nimi), vypocitat jeden Riemannovsky integral (jednoduchy, ale po tech 2 letech to teda byla doba:)  Majerech je velmi hodny a VELMI trpelivy, snazi se poradit - integraly nakonec vicemene v pohode.
{Majerech} losuje - mela jsem stesti a vyfasovala jsem integral. U integralu jsem zmastila, co jsem mohla, jak jsem byla nervozni, takze v ramci vysvetlovani vety o substituci jsem nakonec pocitala i (velmi primitivni) priklad na substituci, ktery jsem nastesti zvladla.
\end{e}
\begin{e}{Report}{0}{Klazar} K Riemannovu integr�lu jsem napsal cca 1,5 A4 definic a v�t, ��dn� d�kaz (n�kter� bych ale v�d�l). To sta�ilo, nedostal jsem snad ani dal�� dopl�uj�c� ot�zky.

\end{e}
\begin{e}{Report}{0}{nezn�m� matematik} ur�it� integr�l, k �emu to je a jak se po��t�, souvislost s derivac�, celkem v pohod�, zn�mku ov�em nev�m

\end{e}
\begin{e}{Report}{0}{Klazar} Riemann�v integr�l. Sta�ilo zaveden�, n�jak� v�ti�ky, vlastnosti, aplikace, drobn� trapas p�i pokusu zintegrovat logaritmus :-).
\end{e}
\begin{e}{Report}{0}{Klazar} uplne vsechno jsem poplet, ale Klazar byl spokojen
\end{e}
\begin{e}{Report}{0}{Matousek} Urcity integral Zdalo se jako snadna otazka, nakonec jsem se i v tomhle dokazal zamotat. Okomentovano slovy: "Takze intergrovat neumite.". Nevzpomnel jsem si akorat na integrovani slozenych funkci(substituce), ostatni jsem po kratsim ci delsim case dal dohromady. Vysledek za 3.
\end{e}
\begin{e}{Report}{0}{Klazar} Riemanuv integral -  Clovek popise papir definicema kolem(deleni, soucty,..), pak padla otazka na primitivni funkci z Riemana, a pry ze staci. Asi 1
\end{e}
\begin{e}{Report}{0}
{Klazar} Riemanuv integral -  Lehk� ot�zka , nicm�n� jsem se tam n�jak zamotal do definic� S a s ( spojil jsem dohromady definici pomoc� suprema/infima toho s/S a definic� pomoc� limity normy d�len� ) , pak cht�l je�t� n�jak� dal�� v�ty , ve kter�ch jsem tak r�zn� pletl p�edpoklady a jednu jsem nedal dohromady v�bec ( respektive jsem m�sto n� dal dohromady asi dv� jin� ) . Tak�e m�j v�kon ��dn� sl�va , nicm�n� Klazar byl fakt hodnej ( to on je v�dycky ) a �ekl , �e m� u� nebude tr�pit a dal mi mysl�m trojku .
\end{e}
\begin{e}{Report}{0}
{Fiala, Kucera} Vicerozmerny integral, Fubiniho veta - tam jsem nemusel dokazovat vubec nic (ani naznakem)
\end{e}
\begin{e}{Report}{0}{Zahradn�k} pocitat jednu nadheru - nejaky viacrozmerny integral s min. 3 typmi substitucii + nejaka ta Fubinka apod.
Tu som ostal trochu zaskoceny, nakolko som cakal zvacsa teoriu, no v prvom priklade bolo teorie minimum, vacsinou som mal dokazovat nejake zaujimave zakonitosti, v druhom pripade zas teoria uplne absentovala. Pokusal som sa spravit mierne teoreticky uvod, no na ten nikto nebol zvedavy. Dalsi mierny problem spocival v tom, ze asi tak v polovicke vypracovania prveho prikladu si sadol ku mne pan Zahradnik so slovami, ze sa pozrie, ako si pocinam, no nejak sa tam zabudol a uz ani nevstal len na konci :) Takze spolocnymi silami sme riesili nove a nove otazky (ako napr. odvodenie vzorca pre objem kuzela a tak). Vacsinu by clovek aj vymyslel po urcitej dobe, aj ten integral by sa dal upocitat, no cakacia doba pana Zahradnika velmi neprekracovala 1 sekundu, takze vysledok bol taky, ze mnoho veci si spocital sam a mne nezostalo nic ine, ako mudro prikyvovat, ze ano, pekna praca :) (obcas to fakt malo nadych konzultacie). Inak prijemny interaktivny styl skusania - rychly sled dvojic - otazka, odpoved (nie vzdy moja :)), ktory by sa hodil v menej stresujucej situacii...  
\end{e}



\end{document}
